\chapter*{Abstract}
\addcontentsline{toc}{chapter}{Abstract}

Accurate and rapid prediction of estuarine hydrodynamics is critical for coastal flood risk management, maritime navigation, and real-time operational forecasting. Traditional physics-based numerical solvers, while highly accurate, are computationally prohibitive for large-scale probabilistic scenario assessments due to the strict Courant number stability constraints of high-resolution numerical integration, which force extremely small simulation time steps. To address this profound computational bottleneck, this thesis introduces \textbf{Hydro-PI-GNN}, a novel parametric Physics-Informed Graph Neural Network surrogate model engineered to dramatically accelerate hydrodynamic simulations on complex, unstructured grids. 

The proposed model was developed, trained, and validated on the Gironde estuary---the largest macrotidal estuary in Western Europe. To provide high-fidelity ground truth data, a rigorous Delft3D Flexible Mesh (Delft3D FM) numerical model comprising 36,271 unstructured finite-volume cells was independently constructed and calibrated. To overcome the phenomenon of ``spectral bias'' inherent in standard deep learning architectures---which routinely causes networks to fail in capturing high-frequency tidal oscillations---Hydro-PI-GNN integrates a mathematically tailored dual-bandwidth Random Fourier Feature (RFF) positional embedding. 

Operating directly on the spatial graph derived from the Delft3D FM mesh, the neural network learns to emulate the complete 2D Shallow Water Equations (SWE). Physical consistency is rigidly enforced via an autograd-based composite loss function that evaluates mass conservation and momentum residuals, incorporating crucial Green--Gauss reconstructed bathymetric bed-slope corrections. The surrogate was optimized using a sophisticated multi-stage curriculum learning strategy and rigorously evaluated using an 80/20 temporal train/test split over a 265-hour tidal cycle.

Validation results on the entirely unseen, held-out test dataset demonstrate that Hydro-PI-GNN achieves exceptional predictive accuracy for the full hydrodynamic state (water surface elevation and depth-averaged velocity vectors) across the entire estuarine domain. Crucially, it dynamically conditions these predictions on continuous variations in ocean tidal amplitude and upstream riverine discharges from the Garonne and Dordogne rivers. By successfully executing complex tidal simulations in fractions of a second during inference, Hydro-PI-GNN achieves a radical reduction in computational time compared to the traditional numerical solver. This research establishes a highly scalable, physically consistent, and open-source methodological framework for deploying high-speed digital twins in vulnerable coastal environments.
